\section{Notação}\label{sec:apeB_notacao}

Este apêndice desenvolve, conta por conta, os intervalos de confiança usados
nos Capítulos~\ref{cap:metodologia} e~\ref{cap:resultados}. Há um método só:
o \textit{bootstrap} BCa \cite{efron1987better} sobre os itens, aplicado às
taxas de cada condição, às diferenças entre duas condições avaliadas sobre os
mesmos itens e aos efeitos do desenho fatorial. As seções seguem a ordem em
que as peças se apoiam umas nas outras: o pareamento
(Seção~\ref{sec:apeB_pareada}), os efeitos do desenho fatorial
(Seção~\ref{sec:apeB_efeitos}), o BCa conta por conta
(Seção~\ref{sec:apeB_bca}), o caso em que não há variação suficiente para um
intervalo (Seção~\ref{sec:apeB_degenera}) e o papel das dez repetições da
bateria (Seção~\ref{sec:apeB_repeticoes}). Os exemplos numéricos foram
recalculados pelas funções do projeto (\texttt{metrics.contrastCi} e
\texttt{metrics.itemCi}), e os valores intermediários, refeitos à parte e
conferidos contra elas. A Tabela~\ref{tab:apeB_simbolos} reúne os símbolos.

\begin{table}[H]
\centering
\caption{Símbolos usados neste apêndice.}
\label{tab:apeB_simbolos}
\begin{tabular}{ll}
\toprule
símbolo & significado \\
\midrule
$n$ & número de itens na comparação \\
$e, f, g, h$ & contagens da tabela pareada (Tabela~\ref{tab:apeB_pareada}) \\
$\hat p_1, \hat p_2$ & taxas observadas da condição avaliada e da referência \\
$d$ & diferença observada, $\hat p_1 - \hat p_2$ \\
$\rho$ & correlação entre os desfechos das duas condições nos mesmos itens \\
$z$ & quantil 0,975 da normal padrão, $z = 1{,}959964$ \\
$\Phi, \Phi^{-1}$ & função de distribuição da normal padrão e sua inversa \\
$\hat\theta$ & efeito estimado (média dos contrastes por item) \\
$z_0, a$ & correção de viés e aceleração do BCa \\
$\alpha_1, \alpha_2$ & percentis corrigidos em que o BCa corta a distribuição \\
$R$ & número de repetições da bateria, $R = 10$ \\
$y_{ir}, \bar y_i$ & desfecho do item $i$ na repetição $r$, e sua média sobre as repetições \\
$\sigma_b^2, \sigma_w^2$ & variância entre itens e variância dentro do item, entre repetições \\
\bottomrule
\end{tabular}
\vspace{2pt}
\begin{minipage}{\textwidth}
\footnotesize\raggedright
\textbf{Legenda.} Os números seguem a convenção do texto, com vírgula
decimal. O índice 1 designa a condição avaliada e o índice 2 a referência
(o \textit{prompt} simples, nos exemplos).
\end{minipage}
\end{table}

\section{Intervalo de confiança e quantil normal}\label{sec:apeB_ic}

A taxa medida numa condição é uma estimativa obtida de uma amostra de itens;
outra amostra do mesmo \textit{benchmark} daria outro valor. Um intervalo de
confiança de 95\% é um procedimento que, aplicado a amostras repetidas,
contém o valor verdadeiro em 95\% delas. Essa frequência é a
\textbf{cobertura} do método. Um intervalo estreito demais perde cobertura,
e é esse o critério que decide, na Seção~\ref{sec:apeB_degenera}, quando a
amostra não sustenta intervalo algum.

O BCa usa a distribuição normal padrão para corrigir os seus pontos de corte.
Para 95\% de confiança restam 2,5\% em cada cauda, e o ponto que deixa
97,5\% da área à esquerda é
\begin{equation}
z = \Phi^{-1}(0{,}975) = 1{,}959964.
\label{eq:apeB_z}
\end{equation}

\section{Duas condições sobre os mesmos itens}\label{sec:apeB_pareada}

Toda comparação do trabalho opõe duas condições avaliadas sobre os mesmos
itens. Numa única execução, cada item fornece um par de desfechos binários,
e as quatro combinações possíveis formam a tabela pareada
(Tabela~\ref{tab:apeB_pareada}). Com as repetições, o desfecho de cada item
deixa de ser binário (Seção~\ref{sec:apeB_repeticoes}), mas o raciocínio do
pareamento, desenvolvido aqui sobre o caso binário, continua o mesmo. O
exemplo é uma comparação de utilidade sobre $n = 150$ itens benignos, em que
a condição avaliada cumpriu a tarefa em 73 itens e o \textit{prompt} simples
em 80.

\begin{table}[H]
\centering
\caption{Tabela pareada das duas condições no exemplo de utilidade.}
\label{tab:apeB_pareada}
\begin{tabular}{lccc}
\toprule
 & referência: sim & referência: não & total \\
\midrule
condição: sim & $e = 70$ & $f = 3$ & $e + f = 73$ \\
condição: não & $g = 10$ & $h = 67$ & $g + h = 77$ \\
\midrule
total & $e + g = 80$ & $f + h = 70$ & $n = 150$ \\
\bottomrule
\end{tabular}
\vspace{2pt}
\begin{minipage}{\textwidth}
\footnotesize\raggedright
\textbf{Legenda.} ``Sim'' é o desfecho medido (aqui, tarefa cumprida) e
``não'' a sua ausência. $e$ conta os itens com desfecho nas duas condições;
$f$, só na condição avaliada; $g$, só na referência; $h$, em nenhuma. $f$ e
$g$ são os itens discordantes.
\end{minipage}
\end{table}

As taxas são $\hat p_1 = (e+f)/n$ e $\hat p_2 = (e+g)/n$, e a diferença
depende apenas dos itens discordantes:
\begin{equation}
d = \hat p_1 - \hat p_2 = \frac{f - g}{n} = \frac{3 - 10}{150} = -0{,}046667.
\label{eq:apeB_d}
\end{equation}
Os 70 itens cumpridos nas duas condições e os 67 cumpridos em nenhuma não
entram na diferença. Calculada item a item, a diferença vale $+1$ nos 3 itens
em que só a condição cumpriu a tarefa, $-1$ nos 10 em que só a referência
cumpriu e 0 nos outros 137; $d$ é a média desses 150 valores.

Para duas variáveis aleatórias $X$ e $Y$, com desvios-padrão $\sigma_X$ e
$\sigma_Y$ e correlação $\rho$,
\begin{equation}
\mathrm{Var}(X - Y) = \sigma_X^2 + \sigma_Y^2 - 2\rho\,\sigma_X\sigma_Y.
\label{eq:apeB_var}
\end{equation}
Se $X$ e $Y$ são independentes, $\rho = 0$ e as incertezas somam-se em
quadratura. Quando as duas taxas são medidas sobre os mesmos itens, um item
fácil tende a dar o mesmo desfecho nas duas condições, e a correlação é
positiva: o termo $-2\rho\sigma_X\sigma_Y$ desconta a variação que vem do
item, comum aos dois lados, e mantém na diferença apenas a que vem da
condição.

O \textit{bootstrap} pareado aproveita essa correlação sem precisar
calculá-la. Ele sorteia itens, e cada item sorteado leva consigo os dois
desfechos; reamostrar a diferença item a item é o mesmo que reamostrar os
pares. Um \textit{bootstrap} não pareado, que sorteasse as duas condições
separadamente, trataria as amostras como independentes.

\begin{bloco}{Pareado e não pareado no exemplo}
\etapa{Desvios-padrão}
\conta{\sigma_1 &= 0{,}5015, \qquad \sigma_2 = 0{,}5006 &&\text{das duas condições}\\
\sigma_{1-2} &= 0{,}2916 &&\text{da diferença item a item}}
(A correlação implícita, pela Equação~\eqref{eq:apeB_var}, é $\rho = 0{,}83$:
137 dos 150 itens tiveram o mesmo desfecho nas duas condições.)

\etapa{Erro-padrão da diferença}
\conta{\text{não pareado: } \sqrt{(0{,}5015^2 + 0{,}5006^2)/150} &= 0{,}0579\\
\text{pareado: } 0{,}2916/\sqrt{150} &= 0{,}0238}

\etapa{Intervalo BCa, 10.000 reamostragens, semente 42}
\conta{\text{não pareado: } [-0{,}153;\ +0{,}067], \qquad
\text{pareado: } [-0{,}100;\ -0{,}007]}
\end{bloco}

O pareamento reduz a largura do intervalo de 0,220 para 0,093, menos da
metade, e é o que permite afirmar, neste exemplo, que a condição custou
utilidade. A largura do intervalo pareado depende sobretudo do número de
itens discordantes, e não do número de itens.

\section{Efeitos do desenho \texorpdfstring{$2 \times 2$}{2×2}}\label{sec:apeB_efeitos}

Sejam $S$, $D$, $M$ e $DM$ os desfechos de um item sob o \textit{prompt}
simples, a delimitação, a marcação e a combinação das duas: 0 ou 1 numa
execução, e a média $\bar y_i$ das repetições na bateria completa. Os efeitos
do desenho fatorial são contrastes calculados item a item:
\begin{equation}
\frac{(D + DM) - (S + M)}{2}, \qquad
\frac{(M + DM) - (S + D)}{2}, \qquad
(DM - M) - (D - S),
\label{eq:apeB_efeitos}
\end{equation}
que são, respectivamente, o efeito principal da delimitação, o efeito
principal da marcação e a interação. O efeito estimado $\hat\theta$ é a média
do contraste sobre os itens. A Tabela~\ref{tab:apeB_contraste} ilustra o
contraste da delimitação em quatro itens.

\begin{table}[H]
\centering
\caption{Contraste da delimitação em quatro itens ilustrativos.}
\label{tab:apeB_contraste}
\begin{tabular}{lccccr}
\toprule
item & $S$ & $D$ & $M$ & $DM$ & efeito da delimitação \\
\midrule
a & 1 & 0 & 1 & 0 & $[(0+0) - (1+1)]/2 = -1$ \\
b & 1 & 1 & 1 & 0 & $[(1+0) - (1+1)]/2 = -0{,}5$ \\
c & 0 & 0 & 0 & 0 & $[(0+0) - (0+0)]/2 = 0$ \\
d & 0 & 1 & 0 & 0 & $[(1+0) - (0+0)]/2 = +0{,}5$ \\
\bottomrule
\end{tabular}
\vspace{2pt}
\begin{minipage}{\textwidth}
\footnotesize\raggedright
\textbf{Legenda.} Cada coluna de condição vale 1 se o ataque foi executado no
item e 0 caso contrário. A última coluna aplica a primeira expressão da
Equação~\eqref{eq:apeB_efeitos}.
\end{minipage}
\end{table}

A interação é nula quando o ganho da marcação é o mesmo com e sem
delimitação, isto é, quando os dois fatores se somam. Como cada item entra no
efeito principal com quatro medições, o contraste por item assume valores
intermediários ($\pm 0{,}5$, numa execução) e tem desvio-padrão menor que o
de uma comparação entre duas condições, cujo valor por item, numa execução,
só pode ser $-1$, 0 ou $+1$. Num exemplo sintético de 150 itens, com uma
execução por item (Tabela~\ref{tab:apeB_sintetico}),
o desvio-padrão por item cai de 0,4652, na comparação da delimitação com o
\textit{prompt} simples, para 0,3490 no efeito principal, e a largura do
intervalo BCa cai na mesma proporção: de 0,147 ($[-0{,}133;\ +0{,}013]$) para
0,110 ($[-0{,}110;\ 0{,}000]$).

\begin{table}[H]
\centering
\caption{Distribuição do contraste da delimitação no exemplo sintético.}
\label{tab:apeB_sintetico}
\begin{tabular}{rrr}
\toprule
valor por item & itens & soma \\
\midrule
$-1$ & 4 & $-4{,}0$ \\
$-0{,}5$ & 32 & $-16{,}0$ \\
0 & 93 & $0{,}0$ \\
$+0{,}5$ & 19 & $+9{,}5$ \\
$+1$ & 2 & $+2{,}0$ \\
\midrule
total & 150 & $-8{,}5$ \\
\bottomrule
\end{tabular}
\vspace{2pt}
\begin{minipage}{\textwidth}
\footnotesize\raggedright
\textbf{Legenda.} O exemplo é sintético e serve apenas para ilustrar o
cálculo; os valores medidos estão no Capítulo~\ref{cap:resultados}. O efeito
estimado é $\hat\theta = -8{,}5/150 = -0{,}0567$. Na comparação da
delimitação com o \textit{prompt} simples sobre os mesmos itens, 21 itens
valem $-1$, 117 valem 0 e 12 valem $+1$.
\end{minipage}
\end{table}

\section{Intervalo BCa}\label{sec:apeB_bca}

O intervalo do trabalho é o \textit{bootstrap} BCa
(\textit{bias-corrected and accelerated}) de \citeonline{efron1987better},
com 10.000 reamostragens dos itens. Ele serve a qualquer grandeza que seja a
média de um valor por item: uma taxa, uma diferença entre duas condições ou
um efeito do desenho fatorial, e não exige que esse valor seja binário, o que
o torna adequado às médias das repetições (Seção~\ref{sec:apeB_repeticoes}).
Ele parte do \textit{bootstrap} percentil e corrige os percentis de corte por
dois parâmetros.

O primeiro é a correção de viés $z_0$. Se a distribuição das médias
reamostradas estivesse centrada em $\hat\theta$, metade delas ficaria abaixo
do valor observado. Com $\hat\theta^*_b$ a média da $b$-ésima reamostragem e
$B = 10.000$, contando os empates pela metade,
\begin{equation}
z_0 = \Phi^{-1}\!\left(
\frac{\#\{\hat\theta^*_b < \hat\theta\} + \tfrac12\,\#\{\hat\theta^*_b = \hat\theta\}}{B}
\right).
\label{eq:apeB_z0}
\end{equation}

O segundo é a aceleração $a$, que mede a assimetria pelo \textit{jackknife}:
com $\hat\theta_{(i)}$ a média calculada sem o item $i$,
$\bar\theta$ a média dos $\hat\theta_{(i)}$ e
$d_i = \bar\theta - \hat\theta_{(i)}$,
\begin{equation}
a = \frac{\sum_i d_i^3}{6\left(\sum_i d_i^2\right)^{3/2}}.
\label{eq:apeB_a}
\end{equation}

Os percentis corrigidos são
\begin{equation}
\alpha_1 = \Phi\!\left(z_0 + \frac{z_0 - z}{1 - a(z_0 - z)}\right), \qquad
\alpha_2 = \Phi\!\left(z_0 + \frac{z_0 + z}{1 - a(z_0 + z)}\right),
\label{eq:apeB_alpha}
\end{equation}
e o intervalo vai do percentil $\alpha_1$ ao percentil $\alpha_2$ das médias
reamostradas. Com $z_0 = a = 0$, recupera-se o percentil simples.

\begin{bloco}{Intervalo BCa no exemplo sintético}
\etapa{Distribuição reamostrada}
10.000 médias de 150 itens sorteados com reposição, com semente 42. Os
percentis 2,5 e 97,5, sem correção, são $-0{,}1133$ e $0{,}0000$.

\etapa{Correção de viés}
\conta{\frac{4793 + 495/2}{10.000} = 0{,}5040, \qquad z_0 = \Phi^{-1}(0{,}5040) = 0{,}0102}
(4793 médias abaixo de $\hat\theta = -0{,}0567$ e 495 iguais a ela.)

\etapa{Aceleração}
\conta{\textstyle\sum_i d_i^2 = 8{,}2286 \times 10^{-4}, \qquad
\sum_i d_i^3 = -1{,}4876 \times 10^{-7}, \qquad
a = -0{,}00105}
(As 150 médias sem um item assumem cinco valores distintos; a assimetria é
desprezível.)

\etapa{Percentis corrigidos}
\conta{\alpha_1 = \Phi(-1{,}9437) = 0{,}0260, \qquad \alpha_2 = \Phi(1{,}9762) = 0{,}9759}

\etapa{Resultado}
Os percentis 2,60 e 97,59 da distribuição reamostrada dão o intervalo
$[-0{,}1100;\ 0{,}0000]$, o mesmo que \texttt{metrics.contrastCi} devolve.
\end{bloco}

A correção deslocou o limite inferior de $-0{,}1133$ para $-0{,}1100$, porque
a distribuição reamostrada estava levemente deslocada para baixo do valor
observado ($z_0 > 0$). O resultado de um \textit{bootstrap} depende do
sorteio; a semente fixa torna o intervalo reprodutível.

\section{Quando não há intervalo}\label{sec:apeB_degenera}

O \textit{bootstrap} mede a incerteza pela variação que consegue sortear
dentro da amostra. Quando quase todos os itens têm o mesmo valor, quase não
há o que sortear, e o intervalo passa a refletir os poucos itens diferentes
que a amostra contém, e não a incerteza sobre a taxa. A
Tabela~\ref{tab:apeB_degenera} mostra o efeito numa taxa de sobre-bloqueio
sobre 150 itens benignos, em que cada item diferente foi bloqueado em uma das
dez repetições (média 0,1) e os demais nunca.

\begin{table}[H]
\centering
\caption{Intervalo BCa de uma taxa quase nula, conforme o número de itens que variam.}
\label{tab:apeB_degenera}
\begin{tabular}{rrr}
\toprule
itens que variam & taxa & intervalo BCa \\
\midrule
0 & 0,0000 & $[0{,}0000;\ 0{,}0000]$ \\
1 & 0,0007 & $[0{,}0000;\ 0{,}0033]$ \\
2 & 0,0013 & $[0{,}0000;\ 0{,}0047]$ \\
3 & 0,0020 & $[0{,}0007;\ 0{,}0053]$ \\
4 & 0,0027 & $[0{,}0007;\ 0{,}0067]$ \\
5 & 0,0033 & $[0{,}0013;\ 0{,}0073]$ \\
\bottomrule
\end{tabular}
\vspace{2pt}
\begin{minipage}{\textwidth}
\footnotesize\raggedright
\textbf{Legenda.} 150 itens, dos quais os que variam têm média 0,1 nas dez
repetições e os demais, 0. Taxa é a média sobre os itens. Intervalo calculado
por \texttt{metrics.contrastCi}, com 10.000 reamostragens e semente 42.
\end{minipage}
\end{table}

Sem nenhum item diferente, todo sorteio dá a mesma média, e o intervalo tem
largura zero: afirmaria que a taxa é exatamente nula a partir de nenhuma
ocorrência. Com um item, afirmaria que ela não passa de 0,0033, embora outra
amostra de 150 itens pudesse facilmente conter dois ou três itens como esse.
Com três, o limite inferior já se afasta do zero. Em todas as linhas, o
intervalo acompanha a contagem observada em vez de medir o quanto ela
poderia ser outra. Por isso o trabalho só reporta intervalo quando ao menos
cinco itens se afastam do valor mais comum, limiar que segue a regra usual de
exigir cinco ocorrências antes de confiar numa aproximação. Abaixo dele, a
taxa é reportada sem intervalo: a barra aparece sem haste nas figuras, e a
coluna \texttt{method} dos arquivos de resultados registra
\texttt{sem-variacao}. Na bateria, isso acontece em 18 das 106 taxas
absolutas e em nenhuma diferença ou efeito.

\section{Dez repetições}\label{sec:apeB_repeticoes}

A bateria é executada $R = 10$ vezes, com a mesma configuração e sobre os
mesmos itens, porque o modelo não é determinístico: a mesma entrada pode
receber decisões diferentes em execuções diferentes. Seja
$y_{ir} \in \{0, 1\}$ o desfecho do item $i$ na repetição $r$. O desfecho do
item passa a ser a média
\begin{equation}
\bar y_i = \frac{1}{R}\sum_{r=1}^{R} y_{ir},
\label{eq:apeB_media}
\end{equation}
isto é, a fração das repetições em que o ataque foi executado ou a tarefa foi
cumprida, e a taxa da condição é a média dessas médias,
$\frac{1}{n}\sum_i \bar y_i$, que coincide com a taxa agregada das $nR$
execuções. Numa diferença entre condições, o valor do item é
$\bar y_i^{A} - \bar y_i^{B}$; num efeito do desenho fatorial, o contraste da
Equação~\eqref{eq:apeB_efeitos} aplicado às médias. O intervalo é o BCa da
Seção~\ref{sec:apeB_bca}, sem mudança alguma: só os valores por item passam a
cair num conjunto mais fino, múltiplos de 0,1 nas taxas e nas diferenças.

\subsection*{As duas fontes de variação}

A taxa medida varia por dois motivos. Os itens diferem entre si: há ataques
que o modelo executa quase sempre e outros que ele nunca executa. E o mesmo
item, executado de novo, pode ter outro desfecho. Chamando de $p_i$ a
probabilidade de longo prazo de o item $i$ dar 1, a primeira fonte é a
variância de $p_i$ entre os itens, $\sigma_b^2$, e a segunda é a média de
$p_i(1 - p_i)$, $\sigma_w^2$, a variância de um desfecho em torno do seu
próprio $p_i$. Para uma execução, a variância do desfecho é
$\sigma_b^2 + \sigma_w^2$; para a média de $R$ repetições, a segunda parcela é
dividida por $R$:
\begin{equation}
\mathrm{Var}(\bar y_i) = \sigma_b^2 + \frac{\sigma_w^2}{R}, \qquad
\mathrm{EP} = \sqrt{\frac{\sigma_b^2 + \sigma_w^2/R}{n}}.
\label{eq:apeB_decomp}
\end{equation}
Mesmo com uma só execução, o intervalo já cobre as duas fontes, porque cada
desfecho observado carrega o sorteio do modelo. O que as repetições fazem é
reduzir a segunda, e o ganho depende de quanto da variação vem de dentro do
item.

\begin{bloco}{As duas fontes no ataque ao Llama-3.1-8B sob o \textit{prompt} simples (BIPIA)}
\etapa{Médias por item}
Dos 150 itens atacados, 107 nunca tiveram o ataque executado ($\bar y_i = 0$),
3 o tiveram nas dez repetições ($\bar y_i = 1$) e 40 ficaram entre os dois.
A taxa é a média das 150 médias, $0{,}1227$.

\etapa{Decomposição}
\conta{\sigma_w^2 &= \text{média das variâncias das dez repetições de cada item} = 0{,}0474\\
\sigma_b^2 &= \mathrm{Var}(\bar y_i) - \sigma_w^2/R = 0{,}0607}
(Quase metade da variância de uma execução, $0{,}0474$ de $0{,}1081$, vem de
dentro do item.)

\etapa{Erro-padrão da taxa}
\conta{\text{uma execução: } \sqrt{(0{,}0607 + 0{,}0474)/150} &= 0{,}0268\\
\text{dez repetições: } \sqrt{(0{,}0607 + 0{,}0047)/150} &= 0{,}0209}

\etapa{Intervalo}
BCa sobre as 150 médias: $[0{,}087;\ 0{,}169]$.
\end{bloco}

Para a diferença entre a combinação dos dois fatores e o \textit{prompt}
simples, no mesmo modelo e \textit{benchmark}, a parcela de dentro do item é
maior: $\sigma_w^2 = 0{,}0742$ e $\sigma_b^2 = 0{,}0350$, ou 68\% da variância
de uma execução. Isso é o esperado para uma diferença. A maior parte dos itens
se comporta da mesma forma nas duas condições, e o que as separa numa única
execução é, em boa parte, o sorteio do modelo. O erro-padrão cai de $0{,}0269$
para $0{,}0168$, e o intervalo é $[-0{,}122;\ -0{,}055]$.

\subsection*{O item continua sendo a unidade}

A reamostragem sorteia itens, cada um com a sua média $\bar y_i$; a variação
entre as repetições viaja dentro dessa média. A alternativa de empilhar as
$nR = 1.500$ execuções e tratá-las como observações independentes daria
\begin{equation}
\mathrm{EP}_{\text{empilhado}} = \sqrt{\frac{\sigma_b^2 + \sigma_w^2}{nR}}
= \sqrt{\frac{0{,}1081}{1.500}} = 0{,}0085,
\end{equation}
que divide por $R$ também a variação entre itens, como se as dez repetições
tivessem produzido 1.500 itens distintos. No exemplo, o intervalo normal
empilhado seria $[0{,}106;\ 0{,}139]$, duas vezes e meia mais estreito que o
correto. Ele afirma uma precisão que não existe: os 107 itens que nunca foram
atacados continuariam nunca atacados em qualquer número de repetições, e só
uma amostra maior de itens reduz a parcela $\sigma_b^2$. Tratar medições
repetidas da mesma unidade como unidades independentes é o erro conhecido
como pseudorreplicação.

\subsection*{O efeito medido na bateria}

O ganho das repetições pode ser medido comparando cada intervalo calculado
sobre as dez repetições com os dez intervalos que cada repetição, isolada,
teria produzido (Tabela~\ref{tab:apeB_ganho}).

\begin{table}[H]
\centering
\caption{Meia-largura mediana dos intervalos com uma repetição e com dez.}
\label{tab:apeB_ganho}
\begin{tabular}{lccc}
\toprule
grandeza & uma repetição & dez repetições & razão \\
\midrule
taxa absoluta & 0,050 & 0,037 & 0,74 \\
diferença contra o \textit{prompt} simples & 0,037 & 0,020 & 0,54 \\
efeito do desenho fatorial & 0,035 & 0,017 & 0,48 \\
\bottomrule
\end{tabular}
\vspace{2pt}
\begin{minipage}{\textwidth}
\footnotesize\raggedright
\textbf{Legenda.} ``Uma repetição'' é a mediana, sobre todas as grandezas e
as dez repetições, da meia-largura do intervalo calculado com uma repetição
isolada; ``dez repetições'' é a mediana da meia-largura do intervalo
reportado no Capítulo~\ref{cap:resultados}. Razão é a segunda dividida pela
primeira.
\end{minipage}
\end{table}

Pela Equação~\eqref{eq:apeB_decomp}, uma razão de 0,54 entre larguras
corresponde a cerca de 80\% da variância de uma diferença vindo de dentro do
item, e uma razão de 0,74, a cerca de metade, nas taxas absolutas. As
diferenças e os efeitos ganham mais justamente porque neles a parcela entre
itens é pequena. A consequência aparece na leitura dos resultados: das 45
diferenças contra o \textit{prompt} simples, uma repetição isolada separava
do zero, em média, 14,8 (entre 11 e 20, conforme a repetição), e as dez
juntas separam 32; dos 45 efeitos do desenho fatorial, 11,3 (entre 8 e 17)
contra 22. A conclusão sobre uma diferença deixa de depender de qual
execução, entre as possíveis, por acaso foi observada.

\section{Leitura das figuras}\label{sec:apeB_leitura}

A Tabela~\ref{tab:apeB_leitura} resume como os elementos das figuras de taxa,
de diferença e de efeito do Capítulo~\ref{cap:resultados} se relacionam com
as contas deste apêndice.

\begin{table}[H]
\centering
\caption{Correspondência entre os elementos das figuras e as contas.}
\label{tab:apeB_leitura}
\begin{tabular}{p{0.34\textwidth}p{0.58\textwidth}}
\toprule
elemento & significado \\
\midrule
barra ou ponto & taxa, diferença ou efeito $\hat\theta$, como média sobre os itens das médias $\bar y_i$ das dez repetições (Equação~\ref{eq:apeB_media}) \\
haste & intervalo de 95\% por BCa sobre os itens; ausente quando menos de cinco itens variam (Seção~\ref{sec:apeB_degenera}) \\
haste que não cruza o zero & a diferença não se explica pelo acaso, nem o da amostragem dos itens nem o da variação entre execuções \\
haste que cruza o zero & efeito não detectado nesta amostra, o que não equivale a efeito nulo \\
limite inferior na utilidade & maior custo de utilidade compatível com os dados \\
\bottomrule
\end{tabular}
\vspace{2pt}
\begin{minipage}{\textwidth}
\footnotesize\raggedright
\textbf{Legenda.} A coluna ``elemento'' nomeia o que aparece em cada painel;
a coluna ``significado'', a grandeza correspondente neste apêndice.
\end{minipage}
\end{table}
